\documentclass[12pt,a4paper,english,reqno]{amsart}
\usepackage[a4paper,footskip=1.5em]{geometry}
\usepackage{amsmath,amssymb,amsthm,mathtools}
\usepackage[mathscr]{euscript}
\usepackage[usenames,dvipsnames]{color}
\usepackage{adjustbox,tikz,calc,graphics,babel,standalone}
\usepackage{subcaption}
\usepackage{csquotes,enumerate,verbatim}
\usepackage[final]{microtype}
\usepackage[numbers]{natbib}
\usetikzlibrary{shapes.misc,calc,intersections,patterns,decorations.pathreplacing}
\usepackage{hyperref}
\hypersetup{colorlinks=true,linkcolor=blue,citecolor=blue,pdfpagemode=UseNone,pdfstartview={XYZ null null 1.00}}

\pagestyle{plain}
\linespread{1.2}
\setlength{\parskip}{3pt}

\theoremstyle{plain}
\newtheorem*{theorem*}{Theorem}
\newtheorem{theorem}{Theorem}[section]
\newtheorem{lemma}[theorem]{Lemma}
\newtheorem{claim}[theorem]{Claim}
\newtheorem{proposition}[theorem]{Proposition}
\newtheorem*{claim*}{Claim}
\newtheorem{corollary}[theorem]{Corollary}
\newtheorem{conjecture}[theorem]{Conjecture}
\newtheorem{problem}[theorem]{Problem}
\newtheorem{question}[theorem]{Question}
\theoremstyle{remark}
\newtheorem*{remark}{Remark}
\newtheorem*{definition}{Definition}

\DeclareMathOperator{\Aut}{Aut}
\DeclareMathOperator{\odd}{odd}
\newcommand{\I}{\mathcal{I}}
\newcommand{\A}{\mathcal{A}}
\newcommand{\B}{\mathcal{B}}
\newcommand{\F}{\mathcal{F}}
\newcommand{\PN}{\mathcal{P}_n}
\newcommand{\N}{\mathbb{N}}
\newcommand{\E}{\mathbb{E}}

\let\emptyset\varnothing
\let\eps\varepsilon
\let\originalleft\left
\let\originalright\right
\renewcommand{\left}{\mathopen{}\mathclose\bgroup\originalleft}
\renewcommand{\right}{\aftergroup\egroup\originalright}

\begin{document}

\title{A refined local lemma lower bound for Folkman's theorem}

\author{Christ Xu}
\email{chrxu@umich.edu}

\date{\today}

\subjclass[2020]{Primary 05D10; Secondary 05D40}

\begin{abstract}
Folkman's theorem asserts that for each $k \in \N$, there exists a natural number $n = F(k)$ such that whenever the elements of $[n]$ are two-coloured, there exists a set $A \subset [n]$ of size $k$ with the property that all the sums of the form $\sum_{x \in B} x$, where $B$ is a nonempty subset of $A$, are contained in $[n]$ and have the same colour. Balogh, Eberhard, Narayanan, Treglown and Wagner~\cite{BENTW} proved that $F(k) \ge 2^{2^{k-1}/k}$ for all $k \in \N$. Their proof uses a random colouring based on the $2$-adic valuation, followed by a union bound over all $k$-sets. We refine the local lemma version of this argument by keeping the total-sum constraint in the dependency count. This gives the stronger estimate
\[
F(k) \ge c k^2 2^{2^{k-1}/(k-1)}
\]
for some absolute constant $c>0$ and all sufficiently large $k$. 
\end{abstract}
\maketitle

\section{Introduction}\label{intro}

Schur's theorem, proved in 1916, is one of the central results of Ramsey theory and asserts that whenever the elements of $\N$ are finitely coloured, there exists a monochromatic set of the form $\{x, y, x+y\}$ for some $x, y \in \N$. About fifty years ago, a wide generalisation of Schur's theorem was obtained independently by Folkman, Rado and Sanders, and this generalisation is now commonly referred to as Folkman's theorem; see~\cite{graham}, for example. To state Folkman's theorem, it will be convenient to have some notation. For $n \in \N$, write $[n]$ for the set $\{1, 2, \dots, n\}$, and for a finite set $A \subset \N$, let
\[
S(A) = \left\{ \sum_{x \in B} x : B \subset A \text{ and } B \ne \emptyset \right\}
\]
denote the set of all nonempty finite sums of elements of $A$. In this language, Folkman's theorem states that for all $k, r \in \N$, there exists a natural number $n = F(k,r)$ such that whenever the elements of $[n]$ are $r$-coloured, there exists a set $A \subset [n]$ of size $k$ such that $S(A)$ is a monochromatic subset of $[n]$. In the case $k=2$, this implies Schur's theorem.

In this note we consider the two-colour Folkman numbers, namely $F(k)=F(k,2)$. All logarithms are to the base $2$. In 1989, Erd\H{o}s and Spencer~\cite{erdosspencer} proved that
\begin{equation}\label{es}
F(k) \ge 2^{c k^2/\log k}
\end{equation}
for all $k \in \N$, where $c>0$ is an absolute constant. This was substantially improved by Balogh, Eberhard, Narayanan, Treglown and Wagner~\cite{BENTW}, who proved
\begin{equation}\label{bentw}
F(k) \ge 2^{2^{k-1}/k}
\end{equation}
for all $k \in \N$. Their proof introduced a random colouring based on the $2$-adic valuation: independently colour each odd integer and then alternate colours along each progression $q,2q,4q,\dots$ with $q$ odd. They showed that, under this colouring, every fixed $k$-set $A$ satisfies
\[
\mathbb P(S(A)\text{ is monochromatic})\le 2^{1-2^{k-1}},
\]
and then applied a union bound over all possible $k$-sets.

The union bound discards the locality of the random colouring. The event that $S(A)$ is monochromatic depends only on the random bits indexed by the odd parts of the elements of $S(A)$. Thus two bad events are independent whenever their sets of odd parts are disjoint. The Lov\'{a}sz local lemma therefore allows one to replace the global count of all $k$-sets by a count of those $k$-sets whose subset sums share a fixed odd part.

Our contribution is a sharpened version of this local lemma argument. The main point is that if a fixed odd integer $q$ occurs as the odd part of a subset sum of a candidate set $B$, then some nonempty subset $T\subseteq B$ has sum $q2^j$. Counting both $T$ and its complement in $B$ using their total-sum constraints, rather than merely bounding each remaining element individually, yields a factorial saving in the dependency degree. This saving translates into a polynomial improvement in the lower bound.

\begin{theorem}\label{mainthm}
There is an absolute constant $c>0$ such that, for all sufficiently large $k$,
\begin{equation}\label{mainbound}
F(k) \ge c k^2 2^{2^{k-1}/(k-1)}.
\end{equation}
Consequently,
\begin{equation}\label{logmainbound}
\log F(k) \ge \frac{2^{k-1}}{k-1}+2\log k-O(1).
\end{equation}
\end{theorem}

Compared with~\eqref{bentw}, this improves the exponent from $2^{k-1}/k$ to $2^{k-1}/(k-1)+2\log k-O(1)$. In multiplicative form, the gain over~\eqref{bentw} is
\[
k^2 2^{2^{k-1}/(k(k-1))-O(1)}.
\]
The exponential part of the gain comes from the local lemma, while the additional factor $k^2$ comes from the refined dependency count.

The paper is organised as follows. In Section~\ref{colouring}, we recall the colouring of Balogh--Eberhard--Narayanan--Treglown--Wagner and the associated probability estimate. In Section~\ref{counting}, we prove the refined counting lemma for fixed odd parts. In Section~\ref{pf}, we apply the Lov\'{a}sz local lemma and prove Theorem~\ref{mainthm}. We conclude with some remarks on the limitations of this method in Section~\ref{conc}.

\section{The Balogh--Eberhard--Narayanan--Treglown--Wagner colouring}\label{colouring}

We recall the random colouring introduced in~\cite{BENTW}. Independently colour each odd integer $q\le N$ with a random bit $X_q\in\{0,1\}$. For each $n\in[N]$, write
\[
n=2^{v_2(n)}\odd(n),
\]
where $\odd(n)$ is odd, and define
\[
\chi(n)=X_{\odd(n)}+v_2(n)\pmod 2.
\]
Thus the colours alternate along every geometric progression $q,2q,4q,8q,\dots$ with $q$ odd.

We shall use the following estimate, which is the central probabilistic ingredient of~\cite{BENTW}.

\begin{claim}\label{probclaim}
For every $k$-set $A\subset [N]$ with $S(A)\subset [N]$,
\[
\mathbb P\big(S(A)\text{ is monochromatic}\big)\le 2^{1-2^{k-1}}.
\]
\end{claim}

\begin{proof}
We include the short argument for completeness.

If $|S(A)|\le 2^k-2$, then two distinct nonempty subsets $B_1,B_2\subseteq A$ have the same sum. Removing $B_1\cap B_2$ from both subsets, we may assume that $B_1$ and $B_2$ are disjoint and still have equal positive sums. Let
\[
x=\sum_{a\in B_1}a=\sum_{a\in B_2}a.
\]
Then $x\in S(A)$ and
\[
2x=\sum_{a\in B_1\cup B_2}a\in S(A),
\]
but the colouring satisfies $\chi(2x)\ne \chi(x)$. Hence $S(A)$ cannot be monochromatic, and the probability is zero.

Now suppose $|S(A)|=2^k-1$, so all nonempty subset sums are distinct. For each odd $m\in\N$, define
\[
G_m=\{m,2m,4m,\dots\}\cap [N].
\]
These progressions partition $[N]$. Let $s=\min_{a\in A}v_2(a)$ and write $a=2^s a'$ for each $a\in A$. At least one of the integers $a'$ is odd. Among the $2^k$ subsets of $A$, exactly half yield
\[
\sum_{a\in B}a'\equiv 1\pmod 2,
\]
since this parity map is a nonzero linear functional on the vector space of subsets of $A$ over $\mathbb F_2$. The empty subset is not among these subsets. Therefore there are exactly $2^{k-1}$ nonempty subsets $B\subseteq A$ for which
\[
v_2\left(\sum_{a\in B}a\right)=s.
\]
The corresponding $2^{k-1}$ sums are distinct and, having the same $2$-adic valuation, lie in distinct progressions $G_m$.

Choose one element of $S(A)$ from each progression $G_m$ that meets $S(A)$, and call the resulting set $B_A$. Then $|B_A|\ge 2^{k-1}$. The colours of the elements of $B_A$ are independent fair bits, since they depend on distinct random variables $X_m$. Hence
\[
\mathbb P(B_A\text{ is monochromatic})=2^{1-|B_A|}\le 2^{1-2^{k-1}}.
\]
Since $S(A)$ monochromatic implies $B_A$ monochromatic, the claim follows.
\end{proof}

\section{A refined fixed-odd-part count}\label{counting}

In this section we prove the counting estimate that sharpens the dependency bound. We use the convention that $\binom{u}{v}=0$ whenever $v<0$ or $u<v$.

\begin{lemma}\label{fixedqlemma}
Let
\[
\A=\left\{B\subset\N: |B|=k,\ \sum_{b\in B}b\le N\right\}.
\]
For a finite set $B\subset\N$, write
\[
O(B)=\{\odd(s):s\in S(B)\}.
\]
Then, for every odd integer $q\le N$,
\begin{equation}\label{fixedqbound}
\#\{B\in\A:q\in O(B)\}
\le
\frac{N^{k-1}}{((k-1)!)^2}
\left(\log N+1+2\binom{2k-2}{k-1}\right).
\end{equation}
\end{lemma}

\begin{proof}
Fix an odd integer $q\le N$. If $B\in\A$ and $q\in O(B)$, then there is a nonempty subset $T\subseteq B$ and an integer $j\ge0$ such that
\[
\sum_{t\in T}t=q2^j.
\]
Put $M=q2^j$ and $r=|T|$. Since $B\in\A$, we necessarily have $M\le N$.

For fixed $M$ and $r$, the number of possible $r$-element sets $T$ with sum $M$ is at most
\[
\frac1{r!}\binom{M-1}{r-1}.
\]
Indeed, $\binom{M-1}{r-1}$ counts ordered positive compositions of $M$ into $r$ parts, while every $r$-element set with sum $M$ gives exactly $r!$ such ordered compositions. The remaining $k-r$ elements of $B$ have total sum at most $N-M$. The number of possible remaining sets is at most
\[
\frac1{(k-r)!}\binom{N-M}{k-r},
\]
because $\binom{N-M}{k-r}$ counts ordered positive $(k-r)$-tuples with total sum at most $N-M$, and every $(k-r)$-element set gives $(k-r)!$ such ordered tuples. These estimates may still overcount, since they do not enforce disjointness between $T$ and the remaining elements, but overcounting is harmless.

Therefore
\begin{equation}\label{precount}
\#\{B\in\A:q\in O(B)\}
\le
\sum_{r=1}^{k}
\sum_{\substack{j\ge0\\q2^j\le N}}
\frac1{r!(k-r)!}
\binom{q2^j-1}{r-1}
\binom{N-q2^j}{k-r}.
\end{equation}

We estimate the right-hand side. For $r=1$,
\[
\sum_{\substack{j\ge0\\q2^j\le N}}
\frac1{(k-1)!}\binom{N-q2^j}{k-1}
\le
(\log N+1)\frac{N^{k-1}}{((k-1)!)^2}.
\]
For $r\ge2$, we use
\[
\binom{q2^j-1}{r-1}\le \frac{(q2^j)^{r-1}}{(r-1)!}
\]
and
\[
\binom{N-q2^j}{k-r}\le \frac{N^{k-r}}{(k-r)!}.
\]
Furthermore,
\[
\sum_{\substack{j\ge0\\q2^j\le N}}(q2^j)^{r-1}\le 2N^{r-1},
\]
since the summands form a geometric progression with ratio $2^{r-1}\ge2$ and largest term at most $N^{r-1}$. Thus the contribution of a fixed $r\ge2$ to~\eqref{precount} is at most
\[
\frac{2N^{k-1}}{r!(r-1)!((k-r)!)^2}.
\]
Since
\[
\frac{((k-1)!)^2}{r!(r-1)!((k-r)!)^2}
=
\frac1r\binom{k-1}{r-1}^2
\le
\binom{k-1}{r-1}^2,
\]
we obtain
\[
\sum_{r=2}^k
\frac{2N^{k-1}}{r!(r-1)!((k-r)!)^2}
\le
\frac{2N^{k-1}}{((k-1)!)^2}
\sum_{r=2}^k\binom{k-1}{r-1}^2.
\]
By Vandermonde's identity,
\[
\sum_{r=1}^k\binom{k-1}{r-1}^2=\binom{2k-2}{k-1}.
\]
Combining the estimates for $r=1$ and $r\ge2$ proves~\eqref{fixedqbound}.
\end{proof}

\section{Proof of the main theorem}\label{pf}

We shall use the symmetric form of the Lov\'{a}sz local lemma.

\begin{lemma}[Symmetric Lov\'{a}sz local lemma]\label{lll}
Let $\{E_i\}_{i\in I}$ be a finite family of events. Suppose each event has probability at most $p$, and each event is mutually independent of all but at most $D$ other events. If $ep(D+1)\le1$, then
\[
\mathbb P\left(\bigcap_{i\in I}E_i^c\right)>0.
\]
\end{lemma}

\begin{proof}[Proof of Theorem~\ref{mainthm}]
Let $N$ be an integer to be chosen later, and use the random colouring described in Section~\ref{colouring}. Define
\[
\A=\left\{A\subset\N: |A|=k,\ \sum_{a\in A}a\le N\right\}.
\]
It is enough to consider this family, since $S(A)\subset[N]$ implies $\sum_{a\in A}a\le N$. For each $A\in\A$, define the bad event
\[
E_A=\{S(A)\text{ is monochromatic}\}.
\]
By Claim~\ref{probclaim},
\begin{equation}\label{pbound}
\mathbb P(E_A)\le p:=2^{1-2^{k-1}}.
\end{equation}

For $A\in\A$, recall that
\[
O(A)=\{\odd(s):s\in S(A)\}.
\]
The event $E_A$ depends only on the random variables $\{X_q:q\in O(A)\}$. Since the random variables $X_q$ are mutually independent, $E_A$ is independent of the sigma-algebra generated by all events $E_B$ with $O(A)\cap O(B)=\emptyset$. We therefore define a dependency graph on $\A$ by joining $A$ and $B$ whenever $O(A)\cap O(B)\ne\emptyset$.

Fix $A\in\A$. Since $|S(A)|\le 2^k-1$, we have $|O(A)|\le2^k-1$. By Lemma~\ref{fixedqlemma}, for each fixed $q\in O(A)$,
\[
\#\{B\in\A:q\in O(B)\}
\le
\frac{N^{k-1}}{((k-1)!)^2}
\left(\log N+1+2\binom{2k-2}{k-1}\right).
\]
Thus the maximum degree $D$ of the dependency graph satisfies
\begin{equation}\label{degreebound}
D
\le
2^k
\frac{N^{k-1}}{((k-1)!)^2}
\left(\log N+1+2\binom{2k-2}{k-1}\right).
\end{equation}
For all sufficiently large $k$ in the range of $N$ chosen below, we have $\log N+1\le2^k$, and also
\[
\binom{2k-2}{k-1}\le4^{k-1}.
\]
Therefore, for all sufficiently large $k$,
\begin{equation}\label{degreebound2}
D\le 8^k\frac{N^{k-1}}{((k-1)!)^2}.
\end{equation}

Let $c_0>0$ be a sufficiently small absolute constant, and set
\begin{equation}\label{Nchoice}
N=\left\lfloor c_0 k^2 2^{2^{k-1}/(k-1)}\right\rfloor.
\end{equation}
Then
\[
N^{k-1}\le c_0^{k-1}k^{2k-2}2^{2^{k-1}}.
\]
By the elementary Stirling lower bound
\[
(k-1)!\ge \left(\frac{k-1}{e}\right)^{k-1},
\]
we have
\[
\frac{k^{2k-2}}{((k-1)!)^2}
\le
\left(\frac{ek}{k-1}\right)^{2k-2}
\le C_1^k
\]
for an absolute constant $C_1>0$. Hence~\eqref{degreebound2} gives
\begin{equation}\label{Dfinal}
D\le C_2^k c_0^{k-1}2^{2^{k-1}}
\end{equation}
for an absolute constant $C_2>0$.

Choose $c_0>0$ so small that $C_2c_0<1/4$. Then, for all sufficiently large $k$,
\[
D\le 2^{-k}2^{2^{k-1}}.
\]
Combining this with~\eqref{pbound}, we get
\[
ep(D+1)
\le
e2^{1-2^{k-1}}\left(2^{-k}2^{2^{k-1}}+1\right)<1
\]
for all sufficiently large $k$. The symmetric Lov\'{a}sz local lemma therefore implies that, with positive probability, none of the events $E_A$ occurs. Consequently, there exists a two-colouring of $[N]$ such that no $k$-set $A$ with $S(A)\subset[N]$ has $S(A)$ monochromatic. Hence $F(k)>N$.

From~\eqref{Nchoice}, after decreasing the absolute constant if necessary, this gives
\[
F(k)\ge c k^2 2^{2^{k-1}/(k-1)}
\]
for all sufficiently large $k$. Taking logarithms gives~\eqref{logmainbound}.
\end{proof}

\section{Concluding remarks}\label{conc}

We conclude with a few remarks.

First, the colouring used here is exactly the colouring of Balogh--Eberhard--Narayanan--Treglown--Wagner~\cite{BENTW}. The improvement comes from two sources: replacing the global union bound by a local lemma argument, and then refining the local dependency count by using the total-sum constraint on both sides of a distinguished subset-sum equation. The first step changes the denominator in the main exponential term from $k$ to $k-1$; the second gives the additional polynomial factor $k^2$.

Second, this method appears to have a natural limitation. In the dependency graph used above, two events are adjacent when their finite-sum sets share an odd part. For a fixed odd integer $q$, the condition that $q$ occur as the odd part of a subset sum imposes essentially one additive equation,
\[
\sum_{t\in T}t=q2^j.
\]
Thus one should expect an $N^{k-1}$-scale family of neighbouring $k$-sets. The refined count in Lemma~\ref{fixedqlemma} improves the factorial and polynomial factors in this estimate, but it does not change the power of $N$. For this reason, the symmetric local lemma with this dependency graph should not be expected to improve the denominator beyond $k-1$.

Third, a further improvement would likely require using more than the uniform worst-case probability bound in Claim~\ref{probclaim}. For each $A$, the same proof actually gives
\[
\mathbb P(E_A)\le 2^{1-|O(A)|}.
\]
Generic sets $A$ may have $|O(A)|$ substantially larger than $2^{k-1}$, while sets with near-minimal $|O(A)|$ should be highly structured. An asymmetric or cluster-expansion local lemma, combined with an inverse theorem for sets with few odd parts among their subset sums, might therefore lead to further progress. We do not pursue this here.
specific case?

\section*{Acknowledgements}
The author thanks J\'{o}zsef Balogh, Sean Eberhard, Bhargav Narayanan, Andrew Treglown, and Adam Zsolt Wagner for their foundational work~\cite{BENTW} on which this note builds.

\begin{thebibliography}{9}

\bibitem{AlonSpencer}
N.~Alon and J.~H.~Spencer,
\emph{The Probabilistic Method},
4th~ed., Wiley, 2016.

\bibitem{BENTW}
J.~Balogh, S.~Eberhard, B.~Narayanan, A.~Treglown, and A.~Z.~Wagner,
An improved lower bound for Folkman's theorem,
\emph{Bull. Lond. Math. Soc.} \textbf{49} (2017), 745--747.

\bibitem{erdosspencer}
P.~Erd\H{o}s and J.~Spencer,
Monochromatic sumsets,
\emph{J. Combin. Theory Ser. A} \textbf{50} (1989), 162--163.

\bibitem{graham}
R.~L.~Graham, B.~L.~Rothschild, and J.~H.~Spencer,
\emph{Ramsey Theory},
2nd~ed., Wiley, 1990.

\bibitem{ub}
R.~L.~Graham,
On the growth of a van der Waerden-like function,
\emph{Integers} \textbf{6} (2006), A29.

\end{thebibliography}

\end{document}
